\scanpage{21}
Soit $\alpha_0\in K_1$, je dis que \unsure{l'on a|on a}
\[
\widehat{\Sigma}_{\alpha_0}\cap\widehat{\Sigma}_i^+=\varnothing
\qquad\text{\unsure{i.e|i.e.}}\qquad
\Sigma_{\alpha_0}\cap\Sigma_i=\varnothing
\]
\unsure{(ce|Ce} qui achèvera la démonstration, car on aura
\[
U=X\setminus\bigcup_{\alpha\in K}\widehat{\Sigma}_\alpha^+
\qquad\left\{\begin{aligned}
&K=K_1\cup\{i\}\\
&\deleted{\unsure{\Sigma_\alpha}}
\end{aligned}\right.
\]
comme dans le lemme.

Or \unsure{comme} on a
\[
E_i\subset U_{J,i}=X\setminus\mathop{\unsure{\rlap{\raisebox{.65ex}{$\scriptstyle-$}}\bigcup|\bigcup}}\limits_{\alpha\in K_0}\widehat{\Sigma}_\alpha^+
\]
cela implique que
\[
E_i\subset X\setminus\widehat{\Sigma}_{\alpha_0}^+=\Sigma_{\unsure{\alpha_0|\alpha}}^-
\]
% margin note written sideways (labels between braces)
\margin{\unsure{j.e.|i.e.}\diagram{\begin{tikzcd}[ampersand replacement=\&,column sep=0pt,row sep=0.9ex,execute at end picture={
\draw ([xshift=-4.0mm]scTop.west) .. controls +(0.7mm,-1mm) and +(1mm,2mm) .. ([xshift=-4.1mm]scMid.west) .. controls +(-1mm,-.8mm) and +(1mm,.5mm) .. ([xshift=-5mm]scMid.west) .. controls +(1mm,-.5mm) and +(1mm,.8mm) .. ([xshift=-4.1mm,yshift=-1mm]scMid.west) .. controls +(.5mm,-2mm) and +(-.6mm,1mm) .. ([xshift=-4.6mm]scEnd.west);
\draw ([xshift=-.9mm]scMid.west) .. controls +(-.4mm,-1mm) and +(.9mm,.7mm) .. ([xshift=-1.1mm]scE.west) .. controls +(-1mm,-.4mm) and +(.7mm,.1mm) .. ([xshift=-1.8mm]scE.west) .. controls +(.7mm,-.1mm) and +(.8mm,.5mm) .. ([xshift=-1.1mm,yshift=-.5mm]scE.west) .. controls +(.4mm,-1mm) and +(-.5mm,.6mm) .. ([xshift=-.8mm]scEnd.west);
\node[rotate=90,inner sep=0pt] at ([xshift=-2.7mm,yshift=-1mm]scE.west) {$\scriptscriptstyle\Sigma_i^-\,\unsure{\rotatebox{-90}{$!$}|\text{ou}|-}\,\Sigma_i^+$};
}]
|[alias=scTop,anchor=west]| \scriptstyle\Sigma_i^+\ \text{ou}\ \Sigma_i^-\\
|[anchor=west]| \scriptstyle E_i\\
|[alias=scMid,anchor=west]| \scriptstyle\unsure{\Sigma_{i\alpha}|\Sigma_{i,\alpha}}\\
|[alias=scE,anchor=west]| \scriptstyle E_\alpha\\
|[alias=scEnd,anchor=west]| \scriptstyle\Sigma_{\alpha_0}^+
\end{tikzcd}}}
d'autre part $\alpha_0\in K_1$ \unsure{signifi|signifie}
\[
\Sigma_{\alpha_0}^+\not\subset\Sigma_i^+
\]
ce qui \deleted{\unsure{impli|impl.}} \unsure{signifie} aussi que $\Sigma_i^+\cap\Sigma_\alpha^+=\varnothing$, cqfd.

\uline{Corollaire du lemme} Pour une classe

\noindent\begin{tabular}{@{}|p{.96\linewidth}@{}}
$V$ de $\Pi_I^{\unsure{\circ|0}}$, son expression
\[
V=X\setminus\bigcup_{\alpha\in K}\widehat{\Sigma}_\alpha^+
\]
du lemme précédent est \uline{unique} -- on a
{\small
\[
\left\{\begin{aligned}
&K=\{\alpha\in I\mid\overline V\cap E_\alpha\ne\varnothing\quad\text{i.e.}\quad\deleted{\unsure{V}}E_\alpha\cap\dot V\ne\varnothing\}\\
&\Sigma_\alpha^+\ \text{est l'unique comp. de }X\setminus\Sigma_\alpha\ \text{contenue dans }X\setminus V
\end{aligned}\right.
\]}
\end{tabular}
% archive sheet number: 20
